\id{IRSTI}{}: 27.35.21

{\bfseries DEVELOPMENT OF MATHEMATICAL MODELS OF THE MAIN UNITS OF THE
SULFUR PRODUCTION UNIT BASED ON SYSTEM ANALYSIS}

{\bfseries B.B.
Orazbayev}\fig{i/image1}[]{\bfseries ,
N.N.
Seit}\fig{i/image1}[]\envelope 

\emph{L.N.Gumilyov Eurasian National University, Astana, Kazakhstan}

\envelope Corresponding author:
nurgulnurmash@gmail.com

A methodology has been proposed for constructing a package of models for
technological systems composed of interconnected units, based on diverse
available information. This methodology is demonstrated through the
development of models for the main units of the sulfur production block.
The models describing the relationship between sulfur yield at the
reactor outlets and their input and operating parameters were identified
as a system of multiple regression equations, using aggregated
experimental, statistical, and expert data, along with the sequential
inclusion of regressors approach. Models assessing the sulfur quality at
the Claus reactor outlet were constructed as fuzzy regression equations.
Parametric identification of the regression coefficients was performed
using the least squares method in the REGRESS software package, while
fuzzy regression coefficients were determined using modified least
squares and α-level sets within the framework of fuzzy set theory. As a
result, mathematical models of the thermoreactor and Claus reactor at
the Atyrau refinery's sulfur production unit were developed.

{\bfseries Keywords:} sulfur production unit, mathematical model, system
analysis, fuzzy information, fuzzy model, structural and parametric
identification.

{\bfseries КҮКІРТ ӨНДІРУ БЛОГЫ НЕГІЗГІ АГРЕГАТТАРЫНЫҢ МАТЕМАТИКАЛЫҚ
МОДЕЛЬДЕР КЕШЕНІН ЖҮЙЕЛІК ТАЛДАУ НЕГІЗІНДН ҚҰРУ}

{\bfseries Б.Б. Оразбаев, Н.Н. Сейт\envelope }

\emph{Л.Н.Гумилев атындағы Еуразия ұлттық университеті, Астана,
Қазақстан,}

e-mail:
nurgulnurmash@gmail.com

Өзара байланысқан агрегаттар кешенінен тұратын технологиялық жүйенің
модельдер пакетін әртүрлі қолжетімді ақпараттар негізінде және жүйелік
талдау тәсілдерін қолдану арқылы құру методикасы әзірленіп, оны күкірт
өндіру блогы мысалында жүзеге асыру көрсетілген. Жинақталған және
өңделген эксперименталдық-статистикалық деректер мен эксперттік
ақпараттар негізінде, сондай-ақ регрессорларды тізбектей қосу тәсілінің
идеясын пайдалана отырып, реакторлар шығысындағы күкірт көлемін
сипаттайтын модельдер құрылымы көп регрессиялық теңдеулер жүйесі түрінде
идентификацияланған. Клаус реакторының шығысындағы күкірт сапасын
бағалайтын модельдер айқын емес регрессиялық теңдеулер түрінде жасалған.
Регрессиялық коэффициенттерді параметрлік идентификациялау ең кіші
квадраттар әдісі арқылы REGRESS бағдарламалық пакетінде жүзеге
асырылған, ал айқын емес коэффициенттерді анықтау α-деңгейлі жиындар
негізінде айқын емес жиындар теориясының тәсілдері арқылы орындалған.
Нәтижесінде, Атырау МӨЗ күкірт өндіру блогының термо және Клаус
реакторларының математикалық модельдері құрылған.

{\bfseries Түйін сөздер:} күкірт өндіру блогы, математикалық модель,
жүйелік талдау, айқын емес ақпарат, айқын емес модель, құрылымдық,
параметрлік идентификациялау.

{\bfseries РАЗРАБОТКА МАТЕМАТИЧЕСКИХ МОДЕЛЕЙ ОСНОВНЫХ АГРЕГАТОВ БЛОКА
ПРОИЗВОДСТВА СЕРЫ НА ОСНОВЕ СИСТЕМНОГО АНАЛИЗА}

{\bfseries Б.Б. Оразбаев, Н.Н. Сейт\envelope }

\emph{Евразийский национальный университет им. Л.Н. Гумилева, Астана,
Казахстан,}

e-mail:
nurgulnurmash@gmail.com

Разработана методика построения пакета моделей технологической системы,
состоящей из взаимосвязанных агрегатов, на основе различных доступных
данных с использованием методов системного анализа, и реализована на
примере блока производства серы. На основе агрегированных и обработанных
экспериментально-статистических данных и экспертной информации, а также
с использованием идеи последовательного включения регрессоров, были
идентифицированы модели, описывающие объем серы на выходе реакторов, в
виде системы многорегрессионных уравнений. Модели, оценивающие качество
серы на выходе Клаус-реактора, построены в форме нечетких регрессионных
уравнений. Параметрическая идентификация регрессионных коэффициентов
выполнялась с применением метода наименьших квадратов через программный
пакет REGRESS, а нечеткие коэффициенты определялись с использованием
методов теории нечетких множеств и преобразования на основе α-уровней. В
результате были построены математические модели термо- и Клаус-реакторов
блока производства серы Атырауского НПЗ.

{\bfseries Ключевые слова:} блок производства серы, математическая модель,
системный анализ, нечеткая информация, нечеткая модель, структурная и
параметрическая идентификация.

{\bfseries Introduction.} In oil refining production, for example at the
Atyrau Oil Refinery, the existing sulfur production unit is a complex
facility, that is, a system consisting of interconnected units
influenced simultaneously by many different parameters {[}1{]}. The main
units of the sulfur production block include reactors, namely the Claus
reactor (R-001), thermal (SWS) reactors (R-002, R-003); condensers
(E-001, E-002, E-004); a furnace (F-002); separators (D-001, D-004); and
pumps (B-001, B-002) {[}2{]}. An analysis of approaches to developing
mathematical models of such complex technological facilities has shown
that issues related to constructing mathematical models of complexes
consisting of interconnected technological units under conditions of
insufficient initial information, as well as system modeling of their
operation, remain insufficiently studied {[}3-8{]}.

When developing mathematical models of industrial technological
facilities and systems, problems arise due to the lack of initial
information required for their construction.

Under conditions of uncertainty arising from the randomness of the
initial information, it is recommended to apply probabilistic modeling
or simulation modeling approaches {[}9-12{]}.

However, if uncertainty arises due to the vagueness of the initial
information-which is common in industrial practice-these approaches
cannot be applied. In such cases, statistical data are insufficient or
completely unavailable, and because the information is vague, it is
impossible to collect quantitative statistical data. As a result, the
axioms of probability theory (such as statistical stability of the
object, repetition of experiments under identical conditions, etc.) are
not satisfied, and therefore there is no basis for applying
probabilistic approaches. Conducting active experiments for the purpose
of collecting statistical data is either impossible or economically
inefficient {[}13, 14{]}. Thus, in these situations, the main available
information consists of vague information expressed in the form of
words, statements, opinions, and judgments based on the knowledge,
experience, and intuition of human decision-makers and expert
specialists. In such cases, it is undoubtedly possible to compensate for
the missing part of the information required for model development by
using vague (fuzzy) information. That is, it is necessary to organize
and conduct expert evaluation {[}15-17{]} and to process its results
using methods of fuzzy set theory {[}18{]}. If the indicated sources of
fuzzy information are reliable-namely, if the decision-makers and expert
specialists are competent and the expert evaluation and its results are
processed correctly-then on the basis of such information it is possible
to construct models that take into account all complex interrelations
between various parameters of complex, non-quantifiable industrial
facilities and processes and that most closely correspond to real
production conditions. The obtained fuzzy models, as a rule, describe
real industrial facilities and problems with higher adequacy than models
constructed using traditional approaches (if such models can be
constructed at all). Therefore, the development of mathematical models
of technological systems in industry based on various available types of
information, including fuzzy information, using the methodology of
systems analysis, is a highly relevant task for both science and
industry.

The aim of this work is to develop a methodology for constructing a
complex of mathematical models for technologically complex units with
difficult quantitative characterization, using the main units of the
sulfur production block as an example, based on systems analysis, as
well as to develop models of the main reactors of the sulfur production
block.

{\bfseries Materials and Methods.} It is necessary to formulate a
methodology for constructing a system of mathematical models of
industrial technological facilities that operate under conditions of
insufficient initial information and uncertainty and are difficult to
describe using quantitative data. Using the Claus reactor and thermal
reactors as examples-where part of the available initial information is
characterized by uncertainty-it is required to construct mathematical
models of reactors R-001, R-002, and R-003 for the purpose of their
optimal control, employing statistical data and fuzzy information
{[}19{]}. Let us formulate the problem of developing mathematical models
of the Claus and thermal reactors using statistical methods, expert
evaluation, and approaches based on fuzzy set theory. As a result of
collecting and processing the available initial information, it has been
established that the statistical data required to estimate the sulfur
yield (output) at the outlets of these reactors can be obtained
experimentally. That is, statistical models determining the sulfur
output at the reactor outlets can be constructed in the form of
regression equations. Since quantitative information for assessing the
quality of the produced sulfur is insufficient and cannot be collected,
models for evaluating sulfur quality should be developed based on
additional fuzzy information.

Thus, through an interconnected internal (parametric) and external
(structural) cyclic process, the problem is formulated of identifying
mathematical models of the sulfur production reactors under study in the
form of the following set of regression and fuzzy regression equations:

\[y_{j} = a_{0j} + \sum_{i = 1}^{n}{a_{ij}x_{ij}} + \sum_{i = 1}^{n}{\sum_{k = i}^{n}{a_{ijk}x_{ij}x_{kj}} + ... + ,j = 1,2};\]

\[\widetilde{y_{j}} = {\widetilde{a}}_{0j} + \sum_{i = n + 1}^{m}{{\widetilde{a}}_{ij}x_{ij}} + \sum_{i = n + 1}^{m}{\sum_{k = i}^{m}{{\widetilde{a}}_{ijk}x_{ij}x_{kj}} + ... + ,j = 3,4}.\]

The constructed models should adequately describe the sulfur yield at
the outlets of the thermal reactors and the Claus reactor
(\(y_{j},\ j = 1,2\)) , as well as the fuzzy indicators characterizing

sulfur quality, namely the mass fraction of sulfur and the mass fraction
(\({\widetilde{y}}_{j},\ j = 3,4\)) of water in the

product obtained from the Claus reactor.

Solution approaches and results. The results of this work include a
methodology for constructing mathematical models of technological
facilities based on various types of information using systems analysis,
as well as models of the Claus and thermal reactors. The units of the
sulfur production block are interconnected, and changes in the operating
parameters of one unit lead to changes in the parameters of others,
thereby affecting the sulfur production process. Therefore, to optimize
and control the sulfur production process, a system of models of the
main units of the technological complex is required, constructed on the
basis of systems analysis and taking into account the influence of
technological parameters of each unit on the intermediate and final
products of the entire complex. In order to develop such a system
(package) of models, this work employs systems analysis methods,
probabilistic and experimental-statistical approaches, as well as expert
evaluation and fuzzy set theory methods {[}20{]}.

Let us formulate a methodology for constructing a complex of
mathematical models of complex industrial technological facilities based
on systems analysis. The model of each facility within a technological
system can be developed using various approaches and methods; that is,
for each unit of the sulfur production complex, a set of models may be
obtained, for example deterministic, statistical, fuzzy, or structured
models.

In order to integrate such diverse models into a single complex (model
package) for the purpose of system-level modeling of the technological
complex, a systems analysis of the advantages and disadvantages of each
model to be constructed is required. This involves defining criteria
that evaluate characteristics such as the costs required to develop the
model, its accuracy, adequacy, applicability for the intended purpose,
and so on, as well as establishing principles for combining the
constructed models into a single package.

For this purpose, the various possible types of models for the main
units of the sulfur production block have been systematically analyzed.
Based on the study of the features of the sulfur production complex and
process {[}21{]}, experimental and expert evaluation data, and an
analysis of modeling approaches for these or similar complex units, the
possible model types for the main units of the sulfur production block
have been assessed. The results of this analysis (model evaluation) are
presented in Table 1. The main criteria used to compare the different
types of models for evaluation are: the availability of information
required to construct the model; the cost (complexity) of model
development; the accuracy (adequacy) of the model; the applicability of
the model for the intended purpose (in our case, solving multi-criteria
optimization and control problems under conditions of uncertainty); and
the possibility of integrating the model into a single system (package)
for the purpose of system-level modeling of the entire facility.

Table 1 presents the evaluations of the possible models for the main
units of the sulfur production block, obtained based on the processing
of the results of the systematic analysis. Based on the information
provided in the table, it is possible to select the most effective types
of models for the main units of sulfur production according to the
chosen criteria.

The construction of regression models for the reactors of the sulfur
production block at the Atyrau Oil Refinery is complicated due to the
lack or insufficiency of proper statistical information, the absence of
specialized industrial instruments, and the low reliability of the
available instruments. Therefore, as an effective tool to supplement the
missing information with additional qualitative data-such as the
knowledge, experience, and judgment of expert specialists-it is
recommended to use expert evaluation methods, while the modeling
approach should involve models and composite methods based on fuzzy set
theory.

The accuracy of the constructed fuzzy or composite models is generally
lower compared to other considered models; however, under conditions of
uncertainty, they are fully sufficient for modeling, optimizing, and
controlling sulfur production processes.

In practice, when information is insufficient, it is necessary to use
all available types of data to construct a model. A model built on such
diverse information is called a \emph{composite model}. However,
constructing a composite model can be inefficient, as it requires
organizing and conducting various types of experiments and studies, as
well as preprocessing the collected data.

{\bfseries Table 1 - Results of the Systematic Analysis of Possible Model
Types for the Main Units of the Sulfur Production Block}

%% \begin{longtable}[]{@{}
%%   >{\raggedright\arraybackslash}p{(\linewidth - 12\tabcolsep) * \real{0.0692}}
%%   >{\centering\arraybackslash}p{(\linewidth - 12\tabcolsep) * \real{0.2036}}
%%   >{\centering\arraybackslash}p{(\linewidth - 12\tabcolsep) * \real{0.3926}}
%%   >{\centering\arraybackslash}p{(\linewidth - 12\tabcolsep) * \real{0.0873}}
%%   >{\centering\arraybackslash}p{(\linewidth - 12\tabcolsep) * \real{0.1018}}
%%   >{\centering\arraybackslash}p{(\linewidth - 12\tabcolsep) * \real{0.0707}}
%%   >{\centering\arraybackslash}p{(\linewidth - 12\tabcolsep) * \real{0.0748}}@{}}
%% \toprule\noalign{}
%% \multirow{2}{=}{\begin{minipage}[b]{\linewidth}\raggedright
%% №
%% \end{minipage}} &
%% \multirow{2}{=}{\centering\arraybackslash \begin{minipage}[b]{\linewidth}\centering
%% {\bfseries Main units of the sulfur production block}
%% \end{minipage}} &
%% \multirow{2}{=}{\centering\arraybackslash \begin{minipage}[b]{\linewidth}\centering
%% {\bfseries Criteria}
%% \end{minipage}} &
%% \multicolumn{4}{>{\centering\arraybackslash}p{(\linewidth - 12\tabcolsep) * \real{0.3345} + 6\tabcolsep}@{}}{%
%% \begin{minipage}[b]{\linewidth}\centering
%% {\bfseries Types of Models}
%% \end{minipage}} \\
%% & & & \begin{minipage}[b]{\linewidth}\centering
%% \begin{quote}
%% {\bfseries Deter-ministic}
%% \end{quote}
%% \end{minipage} & \begin{minipage}[b]{\linewidth}\centering
%% \begin{quote}
%% {\bfseries Statistical}
%% \end{quote}
%% \end{minipage} & \begin{minipage}[b]{\linewidth}\centering
%% \begin{quote}
%% {\bfseries Fuzzy}
%% \end{quote}
%% \end{minipage} & \begin{minipage}[b]{\linewidth}\centering
%% \begin{quote}
%% {\bfseries Composed}
%% \end{quote}
%% \end{minipage} \\
%% \midrule\noalign{}
%% \endhead
%% \bottomrule\noalign{}
%% \endlastfoot
%% 1 & 2 & 3 & 4 & 5 & 6 & 7 \\
%% 1.1 & \multirow{6}{=}{\centering\arraybackslash {\bfseries Reactors}
%% 
%% (R-001,
%% 
%% R-002,
%% 
%% R-003)} & Availability of required information & 2.5 & 4.5 & 4.5 &
%% 5.0 \\
%% 1.2 & & Cost of model development & 1.0 & 4.0 & 3.5 & 3.0 \\
%% 1.3 & & Model accuracy & 4.5 & 3.0 & 2.5 & 4.0 \\
%% 1.4 & & Applicability to the need & 3.5 & 4.0 & 3.5 & 5.0 \\
%% 1.5 & & Possibility of integration into one package & 4.0 & 3.5 & 3.5 &
%% 3.5 \\
%% 1.6 & & Model adequacy & 3.0 & 3.5 & 3.5 & 4.0 \\
%% & & Sum of prices & {\bfseries 18.5} & {\bfseries 22.5} & {\bfseries 21.5} &
%% {\bfseries 24.5} \\
%% 2.1 & \multirow{6}{=}{\centering\arraybackslash {\bfseries Capacitors}
%% 
%% (Е-001,
%% 
%% Е-002,
%% 
%% Е-004)} & Availability of required information & 3.5 & 4.5 & 4.0 &
%% 4.0 \\
%% 2.2 & & Cost of model development & 1.5 & 4.0 & 3.5 & 3.0 \\
%% 2.3 & & Model accuracy & 4.5 & 3.5 & 3.5 & 3.5 \\
%% 2.4 & & Applicability to the need & 4.0 & 4.0 & 4.0 & 4.5 \\
%% 2.5 & & Possibility of integration into one package & 4.5 & 4.0 & 4.0 &
%% 4.0 \\
%% 2.6 & & Model adequacy & 3.5 & 3.5 & 3.5 & 3.5 \\
%% & & Sum of prices & {\bfseries 21.5} & {\bfseries 23.5} & {\bfseries 22.5} &
%% {\bfseries 22.5} \\
%% 3.1 & \multirow{5}{=}{\centering\arraybackslash {\bfseries Oven}
%% 
%% (F-002)} & Availability of required information & 4.0 & 5.0 & 4.5 &
%% 4.0 \\
%% 3.2 & & Cost of model development & 3.0 & 5.0 & 4.0 & 4.0 \\
%% 3.3 & & Model accuracy & 4.5 & 4.5 & 3.5 & 4.0 \\
%% 3.4 & & Applicability to the need & 4.0 & 4.0 & 4.0 & 4.0 \\
%% 3.5 & & Possibility of integration into one package & 4.0 & 4.5 & 4.0 &
%% 5.0 \\
%% 3.6 & & Model adequacy & 4.0 & 4.5 & 4.5 & 5.0 \\
%% & & Sum of prices & {\bfseries 23.5} & {\bfseries 27.5} & {\bfseries 24.5} &
%% {\bfseries 25.0} \\
%% 4.1 & \multirow{5}{=}{\centering\arraybackslash {\bfseries Separators}
%% (D-001,
%% 
%% D-004)} & Availability of required information & 4.5 & 5.0 & 4.0 &
%% 4.5 \\
%% 4.2 & & Cost of model development & 3.0 & 5.0 & 4.0 & 2.5 \\
%% 4.3 & & Model accuracy & 4.5 & 4.5 & 2.0 & 4.0 \\
%% 4.4 & & Applicability to the need & 4.0 & 4.5 & 4.5 & 4.5 \\
%% 4.5 & & Possibility of integration into one package & 3.5 & 4.0 & 3.5 &
%% 4.0 \\
%% 4.6 & & Model adequacy & 4.0 & 4.0 & 3.5 & 4.0 \\
%% & & Sum of prices & {\bfseries 23.5} & {\bfseries 27.0} & {\bfseries 21.5} &
%% {\bfseries 23.5} \\
%% 5.1 & \multirow{5}{=}{\centering\arraybackslash {\bfseries Extractors}
%% 
%% (B-001,
%% 
%% B-002)} & Availability of required information & 4.5 & 4.0 & 4.0 &
%% 4.5 \\
%% 5.2 & & Cost of model development & 5.0 & 4.5 & 4.0 & 4.0 \\
%% 5.3 & & Model accuracy & 5.0 & 4.0 & 4.0 & 4.5 \\
%% 5.4 & & Applicability to the need & 4.5 & 4.5 & 4.0 & 4.5 \\
%% 5.5 & & Possibility of integration into one package & 4.5 & 4.0 & 4.0 &
%% 4.0 \\
%% 5.6 & & Model adequacy & 4.5 & 4.0 & 4.0 & 4.5 \\
%% & & Sum of prices & {\bfseries 28.0} & {\bfseries 25.0} & {\bfseries 24.0} &
%% {\bfseries 26.0} \\
%% \end{longtable}

Note: The evaluation (ranking) uses a score scale (1-5), where 1
represents the lowest score and 5 the highest. The scores may be fuzzy,
i.e., expressed as fuzzy numbers.

\emph{Mathematical Models of the Sulfur Production Block Reactors.} As
determined from the results of the systematic analysis above, the models
of the Claus and thermal reactors are constructed based on various types
of information, including aggregated and processed statistical data,
results of expert evaluations processed using fuzzy set theory methods,
and material and heat balance equations.

To identify the mathematical models of the reactors as multiple
structural and fuzzy regression models, the approach of sequentially
adding regressors {[}22{]} is used, while a modified form of the least
squares method is employed for parametric identification {[}23{]}.

Thus, using the mentioned approach of sequentially adding regressors,
the structure of the mathematical models of the reactors was obtained in
the form of the following multiple and fuzzy regression equations:

\(y_{TR} = a_{0} + \sum_{i = 1}^{4}{a_{i}x_{i}} + \sum_{i = 1}^{4}{\sum_{k = 1}^{4}{a_{ik}x_{i}x_{k}}}\)
(1)

\(y_{RC} = a_{0} + \sum_{i = 5}^{7}{a_{i}x_{i}} + \sum_{i = 5}^{7}{\sum_{k = i}^{7}{a_{ik}x_{i}x_{k}}}\)
(2) \emph{\hfill\break
}\({\widetilde{y}}_{S} = {\widetilde{a}}_{0} + \sum_{i = 5}^{7}{{\widetilde{a}}_{i}x_{i}} + \sum_{i = 5}^{7}{\sum_{k = i}^{7}{{\widetilde{a}}_{ik}x_{i}x_{k}}}\)
(3)

\(\ \ \ \ \ \ \ \ {\widetilde{y}}_{W} = {\widetilde{a}}_{0} + \sum_{i = 1}^{7}{{\widetilde{a}}_{i}x_{i}} + \sum_{i = 5}^{7}{\sum_{k = i}^{7}{{\widetilde{a}}_{ik}x_{i}x_{k}}}\)
(4)

where: \(y_{TR}\) - sulfur yield at the outlet of the thermal reactor;
\(y_{RC}\) - sulfur yield at the outlet of the Claus reactor;
qualitative indicators of the produced sulfur: \({\widetilde{y}}_{S}\) -
fuzzy estimate of the sulfur mass fraction (depending on grade, from
99.20\% to 99.98\%); \({\widetilde{y}}_{W}\) - fuzzy estimate of the
mass fraction of water in the product (depending on grade, from 1.0\% to
0.2\%); input and operating parameters of the reactors affecting the
sulfur production process
(\(x_{i},\ i = \overline{1,4};\ x_{i},\ i = \overline{5,7};\)):
\(x_{1}\) - feed to the thermal reactor (26--28 t/h); \(x_{2}\) -
temperature in the thermal reactor (1000--1413\,°C); \(x_{3}\) -
temperature at the outlet of the thermal reactor (180--350\,°C);
\(x_{4}\) \emph{-} flow rate of combustion air in the thermal reactor
(200--700 Nm³/h); \(x_{5}\) - feed to the Claus reactor (6--8 t/h);
\(x_{6}\) - temperature at the inlet of the Claus reactor
(180--290\,°C); \(x_{7}\) - temperature at the outlet of the Claus
reactor (300-345\,°C);
\(a_{0},a_{i},a_{ik},\ i = \overline{1,4},\ k = \overline{i,4};\) and
\({\widetilde{a}}_{0},{\widetilde{a}}_{i},{\widetilde{a}}_{ik},\ i = \overline{5,7},\ k = \overline{i,7};\)
- regression coefficients to be identified (with crisp and fuzzy
notations): constant term affecting the reactor outlet parameters
\({(a}_{0},{\widetilde{a}}_{0})\) ; linear effects (\(x_{i}\)) ;
quadratic and interaction effects (\(x_{ik}\)).

The obtained models (1)-(4) had their regression coefficients identified
using the REGRESS software package based on the modified least squares
method.

Thus, the results of the parametric identification of the models
describing the relationship between the sulfur yield at the reactor
outlets and their input and operating parameters are as follows:

\[y_{TR} = f_{1}(x_{1}{,x}_{2},x_{4}) = 0.686792x_{1} + 0.012480x_{2} + 0.052000x_{3} + 0.260000x_{4} + 0.025917x_{1}^{2} + 0.0000090x_{2}^{2} + 0.0002080x_{3}^{2} + 0.00005200x_{4}^{2} + 0.000471x_{1}x_{2} + 0.000589x_{1}x_{4} + 0.0000250x_{2}x_{4} + 0.000104x_{3}x_{4}\]

\[y_{RC} = f_{2}(x_{5}{,x}_{6},x_{7}) = 0.6712330x_{5} + 0.0171430x_{6} + 0.0129230x_{7} + 0.0788140x_{5}^{2} + 0.000058x_{6}^{2} + 0.00003x_{7}^{2} + 0.000035x_{5}x_{6} + 0.000589x_{5}x_{7} + 0.0000250x_{6}x_{7}\]

The unknown fuzzy regression coefficients
\({\widetilde{\mathbf{a}}}_{\mathbf{0}}\mathbf{,}{\widetilde{\mathbf{a}}}_{\mathbf{i}}\mathbf{,,\ i =}\overline{\mathbf{5,7}}\)
and
\({\widetilde{\mathbf{a}}}_{\mathbf{ik}}\mathbf{,\ i =}\overline{\mathbf{5,7}}\mathbf{,\ k =}\overline{\mathbf{i,7}}\mathbf{;}\)
of models (3) and (4), which describe the sulfur quality at the Claus
reactor outlet, are evaluated using α-level sets with α = 0.5, 0.75, and
1. As an example, we present \({\widetilde{\mathbf{y}}}_{\mathbf{S}}\)
- the identified fuzzy estimate of the sulfur mass fraction in
the product at the Claus reactor outlet:

\[{\widetilde{y}}_{S} = f_{3}(x_{5}{,x}_{6},x_{7}) = (0.5/0.037008 + 0.75/0.037020 + 1/0.037033 + 0.75/0.037045 + 0.5/0.037057)x_{5} + (0.5/0.076885 + 0.75/0.076900 + 1/0.076915 + 0.75/0.076930 + 0.5/0.076945)x_{6} - (0.5/0.169430 + 0.75/0.169455 + 1/0.169475 + 0.75/0.169495 + 0.5/0.169520)x_{7} + (0.5/0.000010{+ 0.75/0.000020 + 1/0.000027 + 0.75/0.000034 + 0.5/0.000044)x}_{5}^{2} + (0.5/0.000093 + 0.75/0.000108 + 1/0.000118 + 0.75/0.000128 + 0.5/0.000143)x_{6}^{2} - (0.5/0.000007 + 0.75/0.000057 + 1/0.000574 + 0.75/0.000157 + 0.5/0.000207)x_{7}^{2} + (0.5/0.000050 + 0.75/0.000070 + 1/0.000080 + 0.75/0.000090 + 0.5/0.00011)x_{5}x_{6} - (0.5/0.000045 + 0.75/0.000065 + 1/0.000075 + 0.75/0.000085 + 0.5/0.000105)x_{5}x_{7} - (0.5/0.000174 + 0.75/0.000199 + 1/0.000209 + 0.75/0.000219 + 0.5/0.000244)x_{6}x_{7}\]

{\bfseries Results and Discussions.} From the results of the systematic
analysis of the possible model types for the main units of the existing
sulfur production block at the Atyrau Oil Refinery, it is evident-from
the sum of the criteria values (18.5, Table 1)-that due to the
complexity of the reactors and the processes occurring within them, the
difficulty of their study, and the impossibility of obtaining
theoretical information and necessary data, constructing deterministic
models for the reactors is practically infeasible or economically
inefficient.

According to the evaluation results, statistical (stochastic) and fuzzy
models for the reactors of the sulfur production block received
relatively high scores (21.5; 22.5). However, based on the sum of the
criteria values, the most effective type of model for the reactors is
the composite model (24.5). In other words, it is necessary to construct
composite models for the reactors.

For the pumps, based on the criteria scores, constructing a
deterministic model is more advantageous (28.0) compared to other model
types, meaning that a deterministic model is the most efficient choice
for them (Table 1). As shown in the table constructed on the basis of
systems analysis, for the condensers, furnace, and separators, building
statistical models is the most optimal solution (23.5; 27.5; 27,
respectively).

In constructing models of the units included in a technological complex,
a decomposition approach is often used, in which the models of each unit
are developed individually, without considering the integration of the
obtained models into a single system. Solving system-level modeling
tasks in this isolated manner does not yield the desired efficiency or
positive outcomes. Since the operation of an individual unit in a
technological system is interconnected with the operations of other
units, modeling and optimizing a single unit alone cannot fully improve
the performance of the entire system. Therefore, to comprehensively
address the issues of modeling and optimizing a technological system, it
is necessary to create a package of interconnected models that accounts
for the relationships between units. By using such a model package, it
is possible to perform system-level modeling of the technological
complex. Systemic modeling of the technological complex allows for the
identification of ``bottlenecks'' that limit its operation, and
addressing these issues can significantly increase the capacity and
efficiency of the technological complex and its processes.

The integration of individual unit models into a single system (package)
must be carried out in accordance with the operation of the
technological complex process. In this case, the output of one model
(the modeling results) serves as the input for another model (the
initial data required for modeling), and this output may in turn be the
input for a previous or another model. For example, in the sulfur
production complex, the output of the E-001 condenser model - that is,
its modeling results - serves as the input for the E-004 condenser
model, and the results of modeling this condenser are then used as the
input for the R-001 Claus reactor. In turn, the outputs of the Claus
reactor models are required as input data for modeling the
parallel-connected R-002 and R-003 reactors and the E-002 condenser.

The model describing the product of the Claus reactor, that is, the
sulfur yield, was identified as a multiple regression model using
experimental--statistical methods, while the models for assessing sulfur
quality were constructed as fuzzy regression equations based on fuzzy
information obtained from experts.

{\bfseries Conclusion.} Based on the results of the systems analysis, a
methodology for constructing mathematical models of technological
facilities using various types of information has been formulated.
According to the proposed methodology, the model of each unit of a
technological complex is selected from among the possible types of
models based on comparison and selection criteria, identifying the most
effective model for that unit. The possible models of the main units of
the sulfur production block were systematically analyzed, and the most
efficient type of model to be constructed for each unit was determined.

Based on the conducted systems analysis and research results, the
structure and parameters of the mathematical models of the reactors in
the sulfur production block at the Atyrau Oil Refinery were identified
using aggregated and processed information.

The model describing the relationship between the sulfur yield at the
outlets of the thermal and Claus reactors and their input and operating
parameters was identified in the form of multiple regression equations.
The qualitative indicators of sulfur: the mass fractions of sulfur and
water in the product at the outlet of the Claus reactor - were
structurally identified as multiple fuzzy regression equations. Thus, in
this work, using the example of the sulfur production block, a
methodology for constructing a package of models for technological
complexes consisting of interconnected units is presented, along with
the results of its implementation.

{\bfseries Литература}

1. Оразбаев Б.Б., Оразбаева К.Н., Жумадиллаева А.К., Шүйтенов Г.Ж.
Зарыпхан М.Т. Разработка системы моделей взаимосвязанных реакторов
блока производства серы Атырауского нефтеперерабатывающего завода //
Вестник Университета Шакарима. -2025. -Т.2(18). -С.75-86.
\href{https://doi.org/10.53360/2788-7995-2025-2(18)-9}{DOI
10.53360/2788-7995-2025-2(18)-9}.

2. Мурзагалиева Т. Технология переработки нефти и газа: монография //
Фолиант. - 2015. -192c. ISBN 978-601-302-124-9.

3. Zhao Zh-W., Wang D-H.
\href{http://www.sciencedirect.com/science/article/pii/S0895717711007655}{Statistical
inference for generalized random coefficient autoregressive
model}//Mathematical and Computer Modelling. -2012. -Vol.56.
-P.152-166.

\href{https://doi.org/10.1016/j.mcm.2011.12.002}{DOI
10.1016/j.mcm.2011.12.002}.

\setcounter{enumi}{3}

1. Нәдіров Н.К., Сармурзина Р.Ғ., Оразбаева К.Н. Мұнай өңдеу, мұнай
химиясы технологиялық кешендерінің математикалық модельдерін
ақпараттың жетіспеушілігі және айқын еместігі жағдайында құру тәсілін
жасақтау// Доклады НАН РК. -2010. -№2. -Б.77-81.

2. \href{http://www.springerlink.com/content/?Author=Yu.+V.+Sharikov}{Sharikov}
U.V.,
\href{http://www.springerlink.com/content/?Author=P.+A.+Petrov}{Petrov}
P.A.
\href{http://www.springerlink.com/content/br0wn2h316r1v65x/}{Universal
model for catalytic reforming} //
\href{http://www.springerlink.com/content/0009-2355/}{Chemical and
Petroleum Engineering}. -2007. -Vol.43. -Р.580-584.
\href{https://doi.org/10.1007/s10556-007-0103-z}{DOI
10.1007/s10556-007-0103-z}.

3. \href{http://www.sciencedirect.com/science/article/pii/S0895717710004474}{Stampar}
S., Sokolic S.,
\href{http://www.sciencedirect.com/science/article/pii/S0895717710004474}{Karer}
G., Znidarsic A.,
\href{http://www.sciencedirect.com/science/article/pii/S0895717710004474}{Skrjanc}
I. Theoretical and fuzzy modelling of a pharmaceutical batch reactor
//
\href{http://www.sciencedirect.com/science/journal/08957177}{Mathematical
and Computer Modelling}. -2011. -Vol.53. -P.637-645.
\href{https://doi.org/10.1016/j.mcm.2010.09.016}{DOI
10.1016/j.mcm.2010.09.016}.

4. Анчита Х. Переработка тяжелой нефти. Реакторы и моделирование
процессов//Пер.с англ. под ред. Глаголевой О.Ф., Винокурова В.А.
Издательство Профессия. -2015. -592 c. ISBN 978-5-91884-068-9.

5. Беккер В.Ф. Моделирование химико-технологических объектов управления:
учебное пособие. -2-е изд. // РИОР. -2019. -142c.
ISBN~978-5-369-01194-2,~978-5-16-006670-7.

10. Гмурман В.Е. Теория вероятностей и математическая статистика:
учебник. -12-е изд.// Издательство Юрайт. -2024. -480 с. ISBN
978-5-534-00859-3.

11. Акопов А.С. Имитационное моделирование: учебник и практикум для
вузов. -2-е изд.// М.: Издательство Юрайт. - 2024. - 426 c.
ISBN~978-5-534-18379-5.~

12. Кобелев Н.Б., Девятков В.В., Половников В.А. Имитационное
моделирование: учебник. -2-е издание// КУРС. -2024. - 353c. ISBN
978-5-907228-65-8.

13. Воробьев С.Н., Гирина Н.В., Лазарев И.В., Осипов Л.А. Статистическое
моделирование информационных систем: учебное пособие//Издательство ГУАП.
-2010. -15 c. ISBN 978-5-8088-0596-5.

13. Афонин В.В., Федосин С.А. Моделирование систем: учебное пособие //
Национальный Открытый Университет Интуит. -2016. - 269 c. ISBN
978-5-9963-0352-6.

14. Сериков Т.П., Оразбаева К.Н. Интенсификация технологических обьектов
нефтепереработки на основе математических методов: монография // Эверо.
-2006. -157 с. ISBN 9965-769-46-X.

15. Sabzi H.Z., King J.Ph., Abudu Sh. Developing an intelligent expert
system for streamflow prediction, integrated in a dynamic decision
support system for managing multiple reservoirs: A case study // Expert
systems with applications. -2017. -Vol.83. -P.145--163.
\href{https://doi.org/10.1016/j.eswa.2017.04.039}{DOI
10.1016/j.eswa.2017.04.039}.

16. Джарратано Д. Экспертные системы: принципы разработки и
программирование: научное издание пер. с англ. -4-е изд.//Вильямс.
-2007. -1152 с. ISBN~978-5-8459-1156-8, 0-534-38447-1.

17. Оразбаев Б.Б., Курмангазиева Л.Т., Коданова Ш.К. Теория и методы
системного анализа// Издательство Академия Естествознания. -2017. -249с.
ISBN~978-5-91327-498-4.

18. Dubois D. The role of fuzzy sets indecision sciences: Old techniques
and new directions // Fuzzy Sets Systems. -2011. -Vol.184. -P.3-28.
\href{https://doi.org/10.1016/j.fss.2011.06.003}{DOI
10.1016/j.fss.2011.06.003}.

19. Orazbayev В.В., Orazbayeva K.N., Utenova B.E. Development of
Mathematical Models and Modeling of Chemical Engineering Systems under
Uncertainty// Theoretical Foundations of Chemical Engineering. -2014.
-Vol.48. -P.138-147.
\href{https://doi.org/10.1134/S0040579514020092}{DOI
10.1134/S0040579514020092}.

20. Калижанова А.У., Картбаев Т.С., Тогжанова К.О. Системный анализ:
учебное пособие // АУЭС. -2020. -126 с. ISBN 978-601-7939-33-5.

21. Жұманов Қ.Б. Күкірт қышқылы өндірісі: оқу құралы // Фолиант. -2019.
-160б. ISBN~978-601-338-012-4.

22. Валеев С.Г. Регрессивное моделирование при обработке наблюдений:
монография // Наука. -1991. -272 c. ISBN: 5-02-014820-2, 9785020148208.

23. Дилигенская А.Н. Идентификация объектов управления: учебное пособие
// \href{https://samgtu.ru/}{Самарский государственный технический
университет}. -2009. -126 с.

{\bfseries References}

1. Orazbaev B.B., Orazbaeva K.N., Zhumadillaeva A.K., Shүjtenov G.Zh.
Zaryphan M.T. Razrabotka sistemy modelej vzaimosvjazannyh reaktorov
bloka proizvodstva sery Atyrauskogo neftepererabatyvajushhego zavoda //
Vestnik Universiteta Shakarima. -2025. -T.2(18). -S.75-86. DOI
10.53360/2788-7995-2025-2(18)-9. {[}in Russian{]}

2. Murzagalieva T. Tehnologija pererabotki nefti i gaza: monografija //
Foliant. - 2015. -192c. ISBN 978-601-302-124-9. {[}in Russian{]}

3. Zhao Zh-W., Wang D-H.
\href{http://www.sciencedirect.com/science/article/pii/S0895717711007655}{Statistical
inference for generalized random coefficient autoregressive
model}//Mathematical and Computer Modelling. -2012. -Vol.56. -P.152-166.
\href{https://doi.org/10.1016/j.mcm.2011.12.002}{DOI
10.1016/j.mcm.2011.12.002}.

4. Nәdіrov N.K., Sarmurzina R.Ғ., Orazbaeva K.N. Mұnaj өңdeu, mұnaj
himijasy tehnologijalyқ keshenderіnің matematikalyқ
model' derіn aқparattyң zhetіspeushіlіgі zhәne ajқyn
emestіgі zhaғdajynda құru tәsіlіn zhasaқtau// Doklady NAN RK. -2010.
-№2. -B.77-81. {[}in Kazakh{]}

5. \href{http://www.springerlink.com/content/?Author=Yu.+V.+Sharikov}{Sharikov}
U.V.,
\href{http://www.springerlink.com/content/?Author=P.+A.+Petrov}{Petrov}
P.A.
\href{http://www.springerlink.com/content/br0wn2h316r1v65x/}{Universal
model for catalytic reforming} //
\href{http://www.springerlink.com/content/0009-2355/}{Chemical and
Petroleum Engineering}. -2007. -Vol.43. -Р.580-584.
\href{https://doi.org/10.1007/s10556-007-0103-z}{DOI
10.1007/s10556-007-0103-z}.

6. \href{http://www.sciencedirect.com/science/article/pii/S0895717710004474}{Stampar}
S., Sokolic S.,
\href{http://www.sciencedirect.com/science/article/pii/S0895717710004474}{Karer}
G., Znidarsic A.,
\href{http://www.sciencedirect.com/science/article/pii/S0895717710004474}{Skrjanc}
I. Theoretical and fuzzy modelling of a pharmaceutical batch reactor //
\href{http://www.sciencedirect.com/science/journal/08957177}{Mathematical
and Computer Modelling}. -2011. -Vol.53. -P.637-645.
\href{https://doi.org/10.1016/j.mcm.2010.09.016}{DOI
10.1016/j.mcm.2010.09.016}.

7. Anchita H. Pererabotka tjazheloj nefti. Reaktory i modelirovanie
processov//Per.s angl. pod red. Glagolevoj O.F., Vinokurova V.A.
Izdatel' stvo Professija. -2015. -592 c. ISBN
978-5-91884-068-9. {[}in Russian{]}

8. Bekker V.F. Modelirovanie himiko-tehnologicheskih ob\#ektov
upravlenija: uchebnoe posobie. -2-e izd. // RIOR. -2019. -142c. ISBN
978-5-369-01194-2, 978-5-16-006670-7. {[}in Russian{]}

9. Gmurman V.E. Teorija verojatnostej i matematicheskaja statistika:
uchebnik. -12-e izd. // Izdatel' stvo Jurajt. -2024. -
480 s. ISBN 978-5-534-00859-3. {[}in Russian{]}

10. Akopov A.S. Imitacionnoe modelirovanie: uchebnik i praktikum dlja
vuzov. -2-e izd. // M.: Izdatel' stvo Jurajt. - 2024. -
426c. ISBN 978-5-534-18379-5. {[}in Russian{]}

11. Kobelev N.B., Devjatkov V.V., Polovnikov V.A. Imitacionnoe
modelirovanie: uchebnik. -2-e izdanie// KURS. -2024. - 353c. ISBN
978-5-907228-65-8. {[}in Russian{]}

12. Vorob' ev S.N., Girina N.V., Lazarev I.V., Osipov
L.A. Statisticheskoe modelirovanie informacionnyh sistem: uchebnoe
posobie//Izdatel' stvo GUAP. -2010. -15 c. ISBN
978-5-8088-0596-5. {[}in Russian{]}

13. Afonin V.V., Fedosin S.A. Modelirovanie sistem: uchebnoe posobie //
Nacional' nyj Otkrytyj UniversitetIntuit. -2016. - 269 c.
ISBN 978-5-9963-0352-6. {[}in Russian{]}

14. Serikov T.P., Orazbaeva K.N. Intensifikacija tehnologicheskih
ob' ektov neftepererabotki na osnove matematicheskih
metodov: monografija // Jevero. -2006. -157 s. ISBN 9965-769-46-X. {[}in
Russian{]}

15. Sabzi H.Z., King J.Ph., Abudu Sh. Developing an intelligent expert
system for streamflow prediction, integrated in a dynamic decision
support system for managing multiple reservoirs: A case study // Expert
systems with applications. -2017. -Vol.83. -P.145--163.
\href{https://doi.org/10.1016/j.eswa.2017.04.039}{DOI
10.1016/j.eswa.2017.04.039}.

16. Dzharratano D. Jekspertnye sistemy: principy razrabotki i
programmirovanie: nauchnoe izdanie per. s angl. -4-e izd.
//Vil' jams. -2007. -1152 s. ISBN 978-5-8459-1156-8,
0-534-38447-1. {[}in Russian{]}

17. Orazbaev B.B., Kurmangazieva L.T., Kodanova Sh.K. Teorija i metody
sistemnogo analiza// Izdatel' stvo Akademija
Estestvoznanija. -2017. -249s. ISBN 978-5-91327-498-4. {[}in Russian{]}

18. Dubois D. The role of fuzzy sets indecision sciences: Old techniques
and new directions // Fuzzy Sets Systems. -2011. -Vol.184. -P.3-28.
\href{https://doi.org/10.1016/j.fss.2011.06.003}{DOI
10.1016/j.fss.2011.06.003}.

19. Orazbayev В.В., Orazbayeva K.N., Utenova B.E. Development of
Mathematical Models and Modeling of Chemical Engineering Systems under
Uncertainty// Theoretical Foundations of Chemical Engineering. -2014.
-Vol.48. -P.138-147.
\href{https://doi.org/10.1134/S0040579514020092}{DOI
10.1134/S0040579514020092}.

20. Kalizhanova A.U., Kartbaev T.S., Togzhanova K.O. Sistemnyj analiz:
uchebnoe posobie // AUJeS. -2020. -126s. ISBN 978-601-7939-33-5. {[}in
Russian{]}

21. Zhұmanov Қ.B. Kүkіrt қyshқyly өndіrіsі: oқu құraly // Foliant.
-2019. -160b. ISBN 978-601-338-012-4. {[}in Kazakh{]}

22. Valeev S.G. Regressivnoe modelirovanie pri obrabotke nabljudenij:
monografija // Nauka. -1991. -272c. ISBN: 5-02-014820-2, 9785020148208.
{[}in Russian{]}

23. Diligenskaja A.N. Identifikacija ob\#ektov upravlenija: uchebnoe
posobie // Samarskij gosudarstvennyj tehnicheskij universitet. -2009.
-126 s. {[}in Russian{]}

\emph{{\bfseries Information about the authors}}

Orazbayev B.B. - doctor of Technical Sciences professor, L.N. Gumilyov
Eurasian National University, Astana, Kazakhstan, e-mail:
batyr\_o@mail.ru;

Seit N.N. - Senior Lecturer, L.N. Gumilyov Eurasian National University,
Astana, Kazakhstan, e-mail:
nurgulnurmash@gmail.com.

\emph{{\bfseries Сведения об авторах}}

Оразбаев Б.Б. - доктор технических наук, профессор, Евразийский
национальный университет имени Л. Н. Гумилёва, Астана, Казахстан,
e-mail: batyr\_o@mail.ru;

Сейт Н.Н. - старший преподаватель, Евразийский Hациональный университет
имени Л. Н. Гумилёва, Астана, Казахстан, e-mail:
nurgulnurmash@gmail.com.